\section{Introduction}

Dimensional models of psychopathology are displacing categorical diagnosis as
the organizing frame for clinical research. Hierarchical and general-factor
accounts \citep{caspi2014pfactor,lahey2017general,kotov2017hitop,ruggero2019hitop},
normative and biotype modelling \citep{marquand2016,wolfers2018,clementz2016biotypes},
and the transdiagnostic ambition of the Research Domain Criteria
\citep{insel2010rdoc} all reduce a patient to a position in a low-dimensional
latent space. That position is then read, plotted, correlated with outcomes, and
increasingly proposed as a target for stratified care.

A dimensional coordinate, however, is an \emph{estimate}, and an estimate without
its uncertainty is a claim without an error bar. In practice the measurement step
that produces the coordinate is treated as a solved preliminary: incomplete
assessments are imputed or dropped, mixed item types are coerced to a common
scale, and each patient is reported as a point. This hides exactly the quantities
a user of the map needs. How sharply is a given patient actually located --- is a
coordinate pinned down by a dozen observations or floated in from the prior on a
single item? Which of the administered instruments earned that precision, and
which contributed almost nothing? If the assessment is incomplete, what is the
single most valuable thing left to measure? And when a coordinate is later fed
into a prognostic model or a stratification rule, does its uncertainty travel with
it, or is a floated-in prior mean treated as though it were pinned down by data?
Point-estimate, imputation-based scoring cannot answer any of these; the questions
are not about the map's content but about the \emph{quality of its measurement}.

The components needed to answer them already exist, each well characterised in
isolation. Item response theory gives the likelihood of a mixed-type response
given a latent position and, through Fisher information, a currency for how much
each item constrains that position
\citep{lord1980irt,embretsonreise2000,samejima1969,reckase2009mirt}. Expected-a-posteriori
scoring and the Woodbury identity make posterior coordinates and their covariance
computable at scale \citep{bockaitkin1981,harville1997matrix}. Computerized
adaptive testing turns information into a sequential item-selection rule
\citep{vanderlinden2010cat}; value-of-information reasoning turns it into a design
criterion \citep{covertthomas2006}; a principled treatment of missingness
\citep{rubin1976missing,littlerubin2019} makes it legitimate to score a patient
from observed cells alone; and archetypal analysis \citep{cutler1994archetypal}
gives an interpretable low-dimensional summary of a continuous coordinate cloud.
What has not been done is to assemble these into a single object and validate them
\emph{together} on a real, pervasively incomplete, mixed-type transdiagnostic
assessment --- and then to read the resulting map not as a set of coordinates but
as an instrument that can be interrogated about its own measurement.

Here we do exactly that. Using one missingness-native Bayesian bifactor /
exploratory structural-equation model, fit to the harmonised baseline assessment
of $N=9{,}013$ patients with bipolar disorder, schizophrenia or major depression
from the three FACE cohorts, we treat the eight-dimensional dimensional map as a
measurement instrument and ask four operational questions of it. \emph{Where is a
patient, and how sharply} (projection)? \emph{Which items earned that precision}
(information)? \emph{How few items suffice, and where is the instrument bank itself
the limit} (adaptive sampling)? \emph{What is the cheapest next measurement, and
which cohorts are under-measured} (value of information)? We then stress-test the
instrument where a measurement reviewer will push hardest: we show by simulation
that its credible intervals are calibrated, and we quantify how faithfully its
five-archetype convex-blend summary reconstructs the full coordinate.

The claim of this paper is therefore methodological, not clinical: \emph{a
transdiagnostic dimensional map is not a fixed set of point-coordinates but a
self-calibrating measurement instrument} --- one that returns, for every patient,
a posterior coordinate and its uncertainty, quantifies how much information each
indicator contributed (Fisher information, not loading), prescribes the cheapest
next measurement that shrinks that uncertainty, and propagates that uncertainty
into every downstream use. The clinical and biological findings that this map
supports --- including the immunometabolic axis that recurs as its own archetype
corner --- are the subject of a companion paper and are deliberately out of scope
here; where a specific patient or a specific axis appears below, it is an
illustration of an instrument property, never a substantive claim about
psychopathology.
